\documentclass[10pt,letterpaper]{article}
\usepackage{cornell-notes}
\usepackage{mathtools}

\title{Chapter 17: Integration of Ordinary Differential Equations}
\author{}
\date{}

\setDocTitle{Chapter 17: Integration of Ordinary Differential Equations}
\setDocSubtitle{Numerical Methods}
\setDocVersion{v1.0}
\setDocStatus{Study Companion}
\setDocOwner{Numerical Methods Cornell Notes}
\setDocAuthor{}
\setDocDate{\today}
\setDocSummary{Cornell-format study and review notes for Chapter 17: Integration of Ordinary Differential Equations.}

\setCornellCollection{Numerical Methods Cornell Notes}
\setCornellUnitType{Chapter}
\setCornellUnitNumber{17}
\setCornellUnitTitle{Integration of Ordinary Differential Equations}
\setCornellCompanionLabel{Study and review companion}

\hypersetup{pdftitle={Chapter 17: Integration of Ordinary Differential Equations},pdfsubject={Cornell notes},pdfauthor={}}

\begin{document}
\maketitle

\section*{Chapter Overview}
This chapter organizes the mathematical and computational ideas behind integration of ordinary differential equations. The notes preserve the numbered section sequence while emphasizing method selection, explicit assumptions, numerical reliability, cost, and verification.

\section*{Learning Objectives}
Explain the governing problem class; compare Runge-Kutta Method, Adaptive Stepsize Control for Runge-Kutta, Richardson Extrapolation and the Bulirsch-Stoer Method, and Second-Order Conservative Equations; design defensible stopping and verification checks; distinguish problem conditioning from algorithmic stability.

\section*{Prerequisites}
Calculus, differential equations, Taylor methods, norms, and linear systems.

\section*{Operational Workflow}
Scale state and tolerances; choose explicit, implicit, extrapolation, multistep, or structure-preserving integration; estimate local error; accept or reject steps; detect events and stiffness; monitor invariants.

\subsection*{Formula and notation anchor}
\[
y'(t)=f(t,y),\qquad y_{n+1}=y_n+h\sum_i b_i k_i
\]
Runge-Kutta stages $k_i$ sample the derivative at method-specific intermediate states; embedded weights estimate local error.

\begin{warnbox}{Numerical warning}
Stability-limited steps, tolerance mis-scaling, missed events, dense-output inconsistency, order reduction, nonlinear-solve failure, and long-time drift of invariants.
\end{warnbox}

\section{17.0 Introduction}
\begin{cornell}
\noterow{What is the central idea?}{An ODE solver advances state while controlling local error, respecting scales, detecting events, and recognizing stiffness or conserved structure.}
\noterow{How should it be used?}{For Introduction, begin from explicit inputs, outputs, preconditions, and representation choices.  Scale state and tolerances; choose explicit, implicit, extrapolation, multistep, or structure-preserving integration; estimate local error; accept or reject steps; detect events and stiffness; monitor invariants.}
\noterow{Cost and reliability}{Analyze arithmetic work, storage, data movement, and convergence separately.  Relevant failure pressures include stability-limited steps, tolerance mis-scaling, missed events, dense-output inconsistency, order reduction, nonlinear-solve failure, and long-time drift of invariants.  Use scaled tolerances and report both success criteria and diagnostic evidence.}
\noterow{Implementation checklist}{\begin{itemize}
\item Validate dimensions, domains, ordering, and structural assumptions.
\item Guard small divisors, invalid states, overflow, underflow, and nonfinite values.
\item Make mutation, workspace, indexing, and termination behavior explicit.
\item Test a simple exact case, a difficult valid case, a boundary case, and an expected failure.
\end{itemize}}
\end{cornell}
\begin{summarybox}{Section summary}
An ODE solver advances state while controlling local error, respecting scales, detecting events, and recognizing stiffness or conserved structure.  Select this technique only when its assumptions match the data, and accept a result only after an independent residual, error estimate, invariant, or refinement check.
\end{summarybox}

\section{17.1 Runge-Kutta Method}
\begin{cornell}
\noterow{What is the central idea?}{Runge-Kutta schemes combine several within-step derivative samples to obtain a higher-order state update without storing a long history.}
\noterow{How should it be used?}{For Runge-Kutta Method, begin from explicit inputs, outputs, preconditions, and representation choices.  Scale state and tolerances; choose explicit, implicit, extrapolation, multistep, or structure-preserving integration; estimate local error; accept or reject steps; detect events and stiffness; monitor invariants.}
\noterow{Quantitative anchor}{\[
y_{n+1}=y_n+h\sum_i b_i k_i
\]
State the dimensions, scaling, and validity assumptions before using this relation.  Check the resulting residual or invariant rather than trusting an iteration count alone.}
\noterow{Cost and reliability}{Analyze arithmetic work, storage, data movement, and convergence separately.  Relevant failure pressures include stability-limited steps, tolerance mis-scaling, missed events, dense-output inconsistency, order reduction, nonlinear-solve failure, and long-time drift of invariants.  Use scaled tolerances and report both success criteria and diagnostic evidence.}
\noterow{Implementation checklist}{\begin{itemize}
\item Validate dimensions, domains, ordering, and structural assumptions.
\item Guard small divisors, invalid states, overflow, underflow, and nonfinite values.
\item Make mutation, workspace, indexing, and termination behavior explicit.
\item Test a simple exact case, a difficult valid case, a boundary case, and an expected failure.
\end{itemize}}
\end{cornell}
\begin{summarybox}{Section summary}
Runge-Kutta schemes combine several within-step derivative samples to obtain a higher-order state update without storing a long history.  Select this technique only when its assumptions match the data, and accept a result only after an independent residual, error estimate, invariant, or refinement check.
\end{summarybox}

\section{17.2 Adaptive Stepsize Control for Runge-Kutta}
\begin{cornell}
\noterow{What is the central idea?}{Embedded formulas estimate local error from shared stages and adjust the next step according to tolerance and method order.}
\noterow{How should it be used?}{For Adaptive Stepsize Control for Runge-Kutta, begin from explicit inputs, outputs, preconditions, and representation choices.  Scale state and tolerances; choose explicit, implicit, extrapolation, multistep, or structure-preserving integration; estimate local error; accept or reject steps; detect events and stiffness; monitor invariants.}
\noterow{Quantitative anchor}{\[
y_{n+1}=y_n+h\sum_i b_i k_i
\]
State the dimensions, scaling, and validity assumptions before using this relation.  Check the resulting residual or invariant rather than trusting an iteration count alone.}
\noterow{Cost and reliability}{Analyze arithmetic work, storage, data movement, and convergence separately.  Relevant failure pressures include stability-limited steps, tolerance mis-scaling, missed events, dense-output inconsistency, order reduction, nonlinear-solve failure, and long-time drift of invariants.  Use scaled tolerances and report both success criteria and diagnostic evidence.}
\noterow{Implementation checklist}{\begin{itemize}
\item Validate dimensions, domains, ordering, and structural assumptions.
\item Guard small divisors, invalid states, overflow, underflow, and nonfinite values.
\item Make mutation, workspace, indexing, and termination behavior explicit.
\item Test a simple exact case, a difficult valid case, a boundary case, and an expected failure.
\end{itemize}}
\end{cornell}
\begin{summarybox}{Section summary}
Embedded formulas estimate local error from shared stages and adjust the next step according to tolerance and method order.  Select this technique only when its assumptions match the data, and accept a result only after an independent residual, error estimate, invariant, or refinement check.
\end{summarybox}

\section{17.3 Richardson Extrapolation and the Bulirsch-Stoer Method}
\begin{cornell}
\noterow{What is the central idea?}{Extrapolation integrates with a sequence of substeps and removes leading error terms to achieve high accuracy on smooth problems.}
\noterow{How should it be used?}{For Richardson Extrapolation and the Bulirsch-Stoer Method, begin from explicit inputs, outputs, preconditions, and representation choices.  Scale state and tolerances; choose explicit, implicit, extrapolation, multistep, or structure-preserving integration; estimate local error; accept or reject steps; detect events and stiffness; monitor invariants.}
\noterow{Quantitative lens}{Identify the governing equation, recurrence, probability law, or geometric predicate for Richardson Extrapolation and the Bulirsch-Stoer Method.  Define every variable and scale; then derive a residual, error estimate, or invariant that can be checked independently.}
\noterow{Cost and reliability}{Analyze arithmetic work, storage, data movement, and convergence separately.  Relevant failure pressures include stability-limited steps, tolerance mis-scaling, missed events, dense-output inconsistency, order reduction, nonlinear-solve failure, and long-time drift of invariants.  Use scaled tolerances and report both success criteria and diagnostic evidence.}
\noterow{Implementation checklist}{\begin{itemize}
\item Validate dimensions, domains, ordering, and structural assumptions.
\item Guard small divisors, invalid states, overflow, underflow, and nonfinite values.
\item Make mutation, workspace, indexing, and termination behavior explicit.
\item Test a simple exact case, a difficult valid case, a boundary case, and an expected failure.
\end{itemize}}
\end{cornell}
\begin{summarybox}{Section summary}
Extrapolation integrates with a sequence of substeps and removes leading error terms to achieve high accuracy on smooth problems.  Select this technique only when its assumptions match the data, and accept a result only after an independent residual, error estimate, invariant, or refinement check.
\end{summarybox}

\section{17.4 Second-Order Conservative Equations}
\begin{cornell}
\noterow{What is the central idea?}{Structure-aware integrators for second-order systems can preserve phase-space geometry or invariants better over long simulations.}
\noterow{How should it be used?}{For Second-Order Conservative Equations, begin from explicit inputs, outputs, preconditions, and representation choices.  Scale state and tolerances; choose explicit, implicit, extrapolation, multistep, or structure-preserving integration; estimate local error; accept or reject steps; detect events and stiffness; monitor invariants.}
\noterow{Quantitative lens}{Identify the governing equation, recurrence, probability law, or geometric predicate for Second-Order Conservative Equations.  Define every variable and scale; then derive a residual, error estimate, or invariant that can be checked independently.}
\noterow{Cost and reliability}{Analyze arithmetic work, storage, data movement, and convergence separately.  Relevant failure pressures include stability-limited steps, tolerance mis-scaling, missed events, dense-output inconsistency, order reduction, nonlinear-solve failure, and long-time drift of invariants.  Use scaled tolerances and report both success criteria and diagnostic evidence.}
\noterow{Implementation checklist}{\begin{itemize}
\item Validate dimensions, domains, ordering, and structural assumptions.
\item Guard small divisors, invalid states, overflow, underflow, and nonfinite values.
\item Make mutation, workspace, indexing, and termination behavior explicit.
\item Test a simple exact case, a difficult valid case, a boundary case, and an expected failure.
\end{itemize}}
\end{cornell}
\begin{summarybox}{Section summary}
Structure-aware integrators for second-order systems can preserve phase-space geometry or invariants better over long simulations.  Select this technique only when its assumptions match the data, and accept a result only after an independent residual, error estimate, invariant, or refinement check.
\end{summarybox}

\section{17.5 Stiff Sets of Equations}
\begin{cornell}
\noterow{What is the central idea?}{Stiff systems contain rapidly decaying modes that force explicit methods to take stability-limited steps far smaller than accuracy requires.}
\noterow{How should it be used?}{For Stiff Sets of Equations, begin from explicit inputs, outputs, preconditions, and representation choices.  Scale state and tolerances; choose explicit, implicit, extrapolation, multistep, or structure-preserving integration; estimate local error; accept or reject steps; detect events and stiffness; monitor invariants.}
\noterow{Quantitative lens}{Identify the governing equation, recurrence, probability law, or geometric predicate for Stiff Sets of Equations.  Define every variable and scale; then derive a residual, error estimate, or invariant that can be checked independently.}
\noterow{Cost and reliability}{Analyze arithmetic work, storage, data movement, and convergence separately.  Relevant failure pressures include stability-limited steps, tolerance mis-scaling, missed events, dense-output inconsistency, order reduction, nonlinear-solve failure, and long-time drift of invariants.  Use scaled tolerances and report both success criteria and diagnostic evidence.}
\noterow{Implementation checklist}{\begin{itemize}
\item Validate dimensions, domains, ordering, and structural assumptions.
\item Guard small divisors, invalid states, overflow, underflow, and nonfinite values.
\item Make mutation, workspace, indexing, and termination behavior explicit.
\item Test a simple exact case, a difficult valid case, a boundary case, and an expected failure.
\end{itemize}}
\end{cornell}
\begin{summarybox}{Section summary}
Stiff systems contain rapidly decaying modes that force explicit methods to take stability-limited steps far smaller than accuracy requires.  Select this technique only when its assumptions match the data, and accept a result only after an independent residual, error estimate, invariant, or refinement check.
\end{summarybox}

\section{17.6 Multistep, Multivalue, and Predictor-Corrector Methods}
\begin{cornell}
\noterow{What is the central idea?}{History-based methods reuse earlier derivatives for efficiency but require startup, order control, and stability-aware step changes.}
\noterow{How should it be used?}{For Multistep, Multivalue, and Predictor-Corrector Methods, begin from explicit inputs, outputs, preconditions, and representation choices.  Scale state and tolerances; choose explicit, implicit, extrapolation, multistep, or structure-preserving integration; estimate local error; accept or reject steps; detect events and stiffness; monitor invariants.}
\noterow{Quantitative lens}{Identify the governing equation, recurrence, probability law, or geometric predicate for Multistep, Multivalue, and Predictor-Corrector Methods.  Define every variable and scale; then derive a residual, error estimate, or invariant that can be checked independently.}
\noterow{Cost and reliability}{Analyze arithmetic work, storage, data movement, and convergence separately.  Relevant failure pressures include stability-limited steps, tolerance mis-scaling, missed events, dense-output inconsistency, order reduction, nonlinear-solve failure, and long-time drift of invariants.  Use scaled tolerances and report both success criteria and diagnostic evidence.}
\noterow{Implementation checklist}{\begin{itemize}
\item Validate dimensions, domains, ordering, and structural assumptions.
\item Guard small divisors, invalid states, overflow, underflow, and nonfinite values.
\item Make mutation, workspace, indexing, and termination behavior explicit.
\item Test a simple exact case, a difficult valid case, a boundary case, and an expected failure.
\end{itemize}}
\end{cornell}
\begin{summarybox}{Section summary}
History-based methods reuse earlier derivatives for efficiency but require startup, order control, and stability-aware step changes.  Select this technique only when its assumptions match the data, and accept a result only after an independent residual, error estimate, invariant, or refinement check.
\end{summarybox}

\section{17.7 Stochastic Simulation of Chemical Reaction Networks}
\begin{cornell}
\noterow{What is the central idea?}{Event-driven stochastic simulation samples both the next reaction time and reaction channel from state-dependent propensities.}
\noterow{How should it be used?}{For Stochastic Simulation of Chemical Reaction Networks, begin from explicit inputs, outputs, preconditions, and representation choices.  Scale state and tolerances; choose explicit, implicit, extrapolation, multistep, or structure-preserving integration; estimate local error; accept or reject steps; detect events and stiffness; monitor invariants.}
\noterow{Quantitative lens}{Identify the governing equation, recurrence, probability law, or geometric predicate for Stochastic Simulation of Chemical Reaction Networks.  Define every variable and scale; then derive a residual, error estimate, or invariant that can be checked independently.}
\noterow{Cost and reliability}{Analyze arithmetic work, storage, data movement, and convergence separately.  Relevant failure pressures include stability-limited steps, tolerance mis-scaling, missed events, dense-output inconsistency, order reduction, nonlinear-solve failure, and long-time drift of invariants.  Use scaled tolerances and report both success criteria and diagnostic evidence.}
\noterow{Implementation checklist}{\begin{itemize}
\item Validate dimensions, domains, ordering, and structural assumptions.
\item Guard small divisors, invalid states, overflow, underflow, and nonfinite values.
\item Make mutation, workspace, indexing, and termination behavior explicit.
\item Test a simple exact case, a difficult valid case, a boundary case, and an expected failure.
\end{itemize}}
\end{cornell}
\begin{summarybox}{Section summary}
Event-driven stochastic simulation samples both the next reaction time and reaction channel from state-dependent propensities.  Select this technique only when its assumptions match the data, and accept a result only after an independent residual, error estimate, invariant, or refinement check.
\end{summarybox}

\section*{Method comparison}
\begin{center}
\small
\begin{tabularx}{\linewidth}{>{\bfseries\raggedright\arraybackslash}p{0.22\linewidth}XX}
\toprule
Method or lens & Best use & Main caution \\
\midrule
Explicit Runge-Kutta & Nonstiff general problems & Simple adaptive control \\
Implicit method & Stiff problems & Linear or nonlinear solves per step \\
Symplectic method & Long conservative dynamics & Preserves geometric structure \\
\bottomrule
\end{tabularx}
\end{center}

\section*{Worked example}
\begin{exambox}{Independent worked example}
For $y'=y$, $y(0)=1$, one explicit Euler step of size $0.1$ gives $y_1=1.1$; two half steps give $1.1025$, closer to $e^{0.1}\approx1.10517$.  The refinement difference estimates error, but Euler may still be unstable on a stiff decay problem.
\end{exambox}

\section*{Chapter synthesis}
\begin{summarybox}{Chapter synthesis}
The chapter's methods should be viewed as a decision system rather than a list of routines.  Begin with the mathematical problem and its conditioning, exploit structure, select an algorithm with compatible stability and cost, and verify the computed answer with evidence that is independent of the internal iteration.  The recurring implementation discipline is: scale state and tolerances; choose explicit, implicit, extrapolation, multistep, or structure-preserving integration; estimate local error; accept or reject steps; detect events and stiffness; monitor invariants.
\end{summarybox}

\subsection*{Common mistakes and numerical hazards}
\begin{warnbox}{Common mistakes and numerical hazards}
Stability-limited steps, tolerance mis-scaling, missed events, dense-output inconsistency, order reduction, nonlinear-solve failure, and long-time drift of invariants.  A second recurring mistake is reporting only a computed value: always retain tolerances, iteration counts, residuals, refinement behavior, and relevant structural checks.
\end{warnbox}

\subsection*{Retrieval practice}
\textbf{Questions}
\begin{enumerate}
\item What problem does Runge-Kutta Method address, and which assumption is easiest to violate?
\item What problem does Adaptive Stepsize Control for Runge-Kutta address, and which assumption is easiest to violate?
\item What problem does Richardson Extrapolation and the Bulirsch-Stoer Method address, and which assumption is easiest to violate?
\item What problem does Multistep, Multivalue, and Predictor-Corrector Methods address, and which assumption is easiest to violate?
\item What problem does Stochastic Simulation of Chemical Reaction Networks address, and which assumption is easiest to violate?
\item How do conditioning, stability, and the final residual play different roles in this chapter?
\end{enumerate}

\textbf{Brief answers}
\begin{enumerate}
\item Runge-Kutta schemes combine several within-step derivative samples to obtain a higher-order state update without storing a long history. Recheck the section's domain, structure, scaling, and failure diagnostics.
\item Embedded formulas estimate local error from shared stages and adjust the next step according to tolerance and method order. Recheck the section's domain, structure, scaling, and failure diagnostics.
\item Extrapolation integrates with a sequence of substeps and removes leading error terms to achieve high accuracy on smooth problems. Recheck the section's domain, structure, scaling, and failure diagnostics.
\item History-based methods reuse earlier derivatives for efficiency but require startup, order control, and stability-aware step changes. Recheck the section's domain, structure, scaling, and failure diagnostics.
\item Event-driven stochastic simulation samples both the next reaction time and reaction channel from state-dependent propensities. Recheck the section's domain, structure, scaling, and failure diagnostics.
\item Conditioning describes sensitivity of the mathematical problem, stability describes error propagation by the algorithm, and the residual measures satisfaction of the computed equations; none is a universal substitute for the others.
\end{enumerate}

\subsection*{Mastery checklist}
\begin{itemize}
\item[$\square$] I can explain and select Runge-Kutta Method without relying on a memorized code listing.
\item[$\square$] I can explain and select Adaptive Stepsize Control for Runge-Kutta without relying on a memorized code listing.
\item[$\square$] I can explain and select Richardson Extrapolation and the Bulirsch-Stoer Method without relying on a memorized code listing.
\item[$\square$] I can explain and select Second-Order Conservative Equations without relying on a memorized code listing.
\item[$\square$] I can explain and select Stiff Sets of Equations without relying on a memorized code listing.
\item[$\square$] I can state a scaled stopping rule and an independent verification check.
\item[$\square$] I can identify at least one conditioning risk and one algorithmic stability risk.
\item[$\square$] I can estimate time, storage, and data-movement costs at the appropriate level.
\end{itemize}

\end{document}

